% Section 2: The Frame as Method
%
% Status: first-pass draft integrating produced-silence distinction. Will be
% reworked.
%
% Citations needed: SAR data flagging biases, reviewer self-report literature
% (Nisbett \& Wilson 1977 plus operational decision-making), disparate-impact in
% BSA/AML monitoring, cross-institution displacement effects in money laundering.

\section{The Frame as Method}
\label{sec:frame-as-method}

The empty-chair frame's value depends on its generativity. If naming absent-party representation as the organizing principle behind the FS~AI~RMF produces only a more sympathetic vocabulary for compliance theater, the frame has earned nothing. If it produces interpretive leverage on specific controls --- vocabulary for talking about what a control is for, criteria for assessing whether an implementation represents the parties the control was written to protect, diagnostic moves available to examiners reviewing the implementation --- the frame has earned its position as an organizing principle. This section demonstrates the latter through a worked example in BSA/AML transaction monitoring, and in the course of doing so refines the frame itself by distinguishing two kinds of absence.

We choose BSA/AML rather than lending or underwriting deliberately. Lending is where the empty-chair frame is most obvious: the loan applicant is concrete, the harm is concrete, fair-lending controls have decades of regulatory history to draw on, and the absent party's interest is well-rehearsed. The frame applied to lending tends to confirm what readers already believe, which is a poor test of whether the frame is doing work or merely renaming work that was already being done. BSA/AML is harder. The empty chairs are plural and conflicting in ways lending is not. The institution's interest is more aligned with the regulator's than in lending. The customer's interest is ambiguous: protection from fraud and protection from false flagging point in opposite operational directions. If the frame produces clarity here, where the chairs disagree and the controls do not align cleanly with any single absent party, the frame is doing more than renaming.

\subsection{Two kinds of absence}
\label{sec:frame-as-method:two-kinds-of-absence}

Before enumerating the empty chairs in BSA/AML, we refine the frame. The opening movement treated absence as a single category: parties not in the room when a decision is made. In practice, absence has two structurally distinct forms, and the diagnostic moves available to address them differ.

\emph{Structural absence} describes parties whose absence is given by the situation: parties who do not yet exist, who have no embodied agent, or whose absence is a designed feature of the regulatory regime. The future compliance officer, the not-yet-victimized customer, the financial system as such, the criminal we are designing to exclude: these are absent because the architecture of the situation places them outside the decision context. The frame's diagnostic move on structural absence is \emph{how does the architecture make this party audible to the parties who are present?}

\emph{Produced absence} describes parties whose absence is created by the architecture itself, then read by the architecture as evidence of correctness. The customer who would have appealed if appeal were not prohibitively expensive. The borrower who stopped applying after repeated denial and whose absence from the application data is read as evidence of self-selection out of credit demand. The reviewer whose actual reasoning is structurally erased by a documentation regime that captures only post-hoc justification, where the absence of contradicting evidence in the record is read as evidence the documentation is faithful. In each case, the architecture has costs or design choices that suppress a voice, and the architecture then interprets the suppressed voice's silence as endorsement of the decision that suppressed it.

The distinction matters because produced absence is invisible to the empty-chair frame as initially stated. \emph{Make the absent party audible} is the wrong operation when the architecture is producing the silence: the chair appears empty because the architecture removed it, and asking how to give voice to the chair while leaving the removal-mechanism intact will fail. The diagnostic move on produced absence is \emph{what is the architecture doing that creates this silence, and what is it inferring from the silence it created?}

This distinction sharpens position one of the paper. Compliance theater is partly the failure to represent structural absences. But compliance theater's persistence is more substantially about produced absences: architectures that suppress challenge, then read the absence of challenge as evidence the architecture is correct. The dishonest middle is not merely a failure of representation. It is an active production of silence, interpreted as endorsement.

\subsection{The empty chairs in BSA/AML monitoring}
\label{sec:frame-as-method:bsa-aml-empty-chairs}

Six absent parties have interests in how a community bank monitors transactions for suspicious activity \parencite{ffiecBSAAMLManual}. Each is structural, produced, or both. We mark which.

\emph{The next victim of fraud whose pattern this transaction is part of.} Structural absence. When a transaction is part of an emerging fraud pattern, the institution's monitoring decision affects not the customer making the transaction but the customer who will be victimized in the next iteration. The next victim is fully absent from the transaction record and from the institution's customer relationship by the structure of time and the limits of prediction. They are represented in the monitoring architecture or they are not represented at all.

\emph{The customer whose account is frozen on a false positive.} Mixed: structural and produced. The customer is present in the transaction record but absent from the decision context that determined how their pattern was scored; that absence is structural, given that the model evaluates the transaction without consulting the customer. But the customer's silence after the freeze is partly produced: appeal processes are typically opaque, expensive, and slow, and many false-positive customers absorb the cost rather than contest it. The institution's apparent flagging accuracy is sustained in part by appeal-cost dynamics that suppress challenge from the falsely flagged.

\emph{The marginalized customer whose transaction patterns are flagged at elevated rates.} Mixed: structural and produced. Models trained on historical SAR data are at risk of inheriting historical flagging biases, especially where enforcement and reporting have differed across communities even when underlying risk has not \parencite{treasuryDeRisking2023}. Communities whose transaction patterns differ from the model's training distribution (recent immigrants using money-transfer services, cash-economy participants, members of religious or cultural communities with non-standard remittance patterns) can experience flagging rates that do not track underlying risk consistently across populations. The training-distribution gap is structural absence: the model was not trained to represent these populations as separate. But the appeal-cost dynamic is produced absence: the population most likely to be falsely flagged is also the population least equipped to contest the flag, and the institution's stable false-positive rate across populations is partly an artifact of differential appeal costs rather than evidence of unbiased flagging.

\emph{The institution itself, in the form of its future enforcement exposure.} Structural absence. The bank that fails to file required SARs faces regulatory action; the bank that over-files faces examiner skepticism about its risk calibration. Both are forms of enforcement exposure that compound over time. The institution's future self --- the bank in five years, when current monitoring decisions become the subject of an examination --- is an empty chair the present institution can either represent through architectural attestation or fail to represent through reliance on reconstructed reasoning.

\emph{The financial system, in the form of monitoring quality across institutions.} Structural absence. BSA/AML is a coordinated regulatory regime; one institution's monitoring inadequacy creates externalities for institutions whose monitoring is more rigorous, because illicit activity predictably gravitates toward weaker control environments. The system as a whole is an empty chair represented by inter-agency information sharing, FinCEN reporting, and cross-institution examination practice. The community bank's monitoring architecture either contributes to this representation or free-rides on it.

\emph{The criminal whose laundering depends on monitoring inadequacy.} Structural absence by intentional design. We name this empty chair explicitly because its absence in the architecture is intentional and structural, and because conflating ``we should not represent the criminal'' with ``we should not consider the criminal's interest in our threat model'' is a category error. The criminal's interest is the threat the architecture is designed against. The frame requires naming the chair we are designing to exclude. We also note the boundary case: the criminal and the marginalized customer can be hard to distinguish from behavior patterns alone, and architectures that conflate the two by design produce a particularly vicious form of produced absence: the marginalized customer flagged as criminal, whose appeal is expensive and whose silence after flagging is read as confirmation of suspicion.

The 230 control objectives do not generally describe controls in these terms. They specify governance outcomes for AI systems --- inventory attributes, oversight mechanisms, explanation practices, monitoring processes --- without naming every party whose interest an implementation might protect, and without distinguishing structural from produced absence. The empty-chair frame's first contribution is a vocabulary for examining those implementation choices: which chairs a control objective can represent, which chairs a particular implementation actually represents, and whether any remaining absence is given by the setting or produced by the implementation itself.

\subsection{Walking three examples through the frame}
\label{sec:frame-as-method:walking-three-controls}

The examples below apply four cross-sector FS~AI~RMF objectives to a community bank's BSA/AML transaction-monitoring system. They are not BSA/AML-specific requirements. For each example, we distinguish the objective's text from its implementation guidance and example evidence, and both from the bank implementation we ask the reader to consider. This distinction matters: an objective can support several implementations, and an implementation may support inferences the objective neither requires nor warrants. The objective identifiers and descriptions below are from the full Control Objective Reference Guide, version~1.0 \parencite{fssccFSAIRMF2026}.

\emph{Inventory: GV-1.6.1, AI Inventory Management.} The objective requires a regularly updated inventory of AI systems that includes purpose, data sources, responsible personnel, risk ratings, criticality, dependencies, compliance requirements, and decommissioning attributes. Its implementation guidance recommends a standardized template, assigned maintenance responsibility, periodic or trigger-based review, and access controls. Its examples of effective evidence include inventory records, ownership details, change histories, and audit trails.

Applied to transaction monitoring, a current inventory entry establishes something important: the institution has identified the system, assigned responsibility, and recorded specified lifecycle and risk attributes. It allows a future compliance officer or examiner to locate the system in the institution's governance landscape. It does not, by itself, establish that the listed data sources represent affected populations adequately, that the documented purpose matches operational use, or that a particular alert was well founded. Those are different propositions requiring different evidence.

The empty-chair question reveals the difference. The inventory directly represents the institution's future self and the future compliance officer by preserving organizational location and responsibility. It represents a customer affected by an alert only indirectly, through attributes that make later inquiry possible. A bank that treats the existence of a complete inventory entry as evidence that individual outputs are justified has asked the inventory to establish more than it can. Closing that gap requires evidence linked to the deployed system and the decision under review: the operative model version, relevant input lineage, validation results, and the record of how the output entered the decision process. The inventory is the index into that evidence, not a substitute for it.

\emph{Human oversight: GV-3.2.2, Monitoring Human-AI Oversight.} The objective requires mechanisms such as audits, performance evaluations, and feedback loops to assess human--AI interactions, including automation complacency and inappropriate overreliance. Its guidance asks institutions to examine how well operators detect issues, respond to alerts, and understand system limitations. Its examples include interaction audits, operator evaluations, feedback mechanisms, training, and ongoing assessment of oversight quality.

In a transaction-monitoring implementation, completed audits and operator evaluations establish that the institution examined specified aspects of human oversight. Measures of response quality, override patterns, or detected errors can show whether reviewers are engaging with alerts rather than merely accepting them. But a signed review record establishes only that a reviewer performed the recorded act under the institution's procedure. It does not, without further evidence, establish that the alert was substantively correct or that the reviewer independently evaluated every consideration material to the disposition.

This distinction serves both the falsely flagged customer and the next potential victim. For the former, a review signature is not evidence that the alert's basis was sound. For the latter, an override is not evidence that the reviewer noticed every risk the system surfaced. An implementation can narrow the gap by preserving what information was presented, what sources the reviewer consulted, whether the reviewer accepted or changed the recommended disposition, and which uncertainty or evidence triggered that action. Even this contemporaneous record has limits: it documents observable review activity, not unmediated access to the reviewer's mental process. The objective asks whether human oversight is effective; the implementation must specify which evidence makes that effectiveness assessable.

\emph{Explanation and decision context: MS-2.9.1, Model Explainability and Validation, and MS-2.9.2, Interpretability and Decision Context.} MS-2.9.1 calls for structured model explanations covering matters such as design choices, training-data sources, feature importance, and decision pathways, together with validation for stability, fairness, performance, bias, and limitations. Its example controls expressly include interpretable models and post-hoc tools such as SHAP and LIME, with documentation of methods, parameters, assumptions, and limitations. MS-2.9.2 requires outputs and decision rationales to be interpreted for their use context and intended audience. Its guidance and examples emphasize clarity, accessibility, human review, feedback, and testing the usefulness of explanations.

These objectives recognize several legitimate evidentiary functions. Model documentation can establish declared design choices and training sources. Validation can establish measured performance under specified tests. A SHAP or LIME report can characterize aspects of model behavior under the method's construction. User testing can establish whether an intended audience finds an explanation understandable. None of these results automatically establishes the others. In particular, clarity to a reviewer is not evidence that an explanation faithfully represents every computational dependency material to an output, and a technically faithful attribution is not evidence that the attributed model behavior is lawful or policy-conforming.

The empty-chair reading therefore asks which proposition matters to the affected party. A customer challenging a false alert may need to know which evidence supported the institution's action and whether that evidence was applied consistently with policy. An examiner may need to know the model's measured behavior across populations and whether the explanation procedure is stable under relevant conditions. A reviewer may need an account calibrated to the decision they are authorized to make. Treating one explanation artifact as satisfying all three needs produces a gap precisely where the objectives invite contextual implementation. Closing or disclosing that gap requires the institution to name the artifact, the proposition it supports, the method and assumptions that produced it, and the additional evidence needed for any broader inference.

\subsection{What the frame produced}
\label{sec:frame-as-method:what-the-frame-produced}

Four moves, in this section, the control objectives alone did not produce. First, an explicit accounting of which absent parties have interests in BSA/AML monitoring, including the chair we are designing to exclude. Second, a distinction between structural and produced absence, with different diagnostic moves for each. Third, a vocabulary for separating objective text, implementation evidence, and the further inferences drawn from that evidence. Fourth, a diagnostic question pair --- \emph{which empty chair does this implementation represent, and what silence does it produce that it then reads as endorsement?} --- that supplements the compliance question --- \emph{does this implementation satisfy the objective?}

These are not solutions to the problems the objectives address. They are a frame for interpreting the objectives and testing what their implementations establish. The frame's claim to organizing-principle status rests on whether moves of this kind generalize across controls and regulated decision contexts, which the remainder of this paper develops. What this section establishes is narrower: the frame makes visible a distinction among the required governance outcome, the artifact offered as evidence, and the conclusion drawn from that artifact. The structural-versus-produced distinction then asks why any missing evidence is absent and whose interests bear the consequence.
